% Part: normal-modal-logic
% Chapter: syntax-and-semantics
% Section: truth-in-model

\documentclass[../../../include/open-logic-section]{subfiles}

\begin{document}

\olfileid{nml}{syn}{tru}

\olsection{નિદર્શનમાં સત્યતા}

ક્યારેક કોઈ નિદર્શનના દરેક જગતમાં કયાં !!{formula}s સત્ય છે
તે જાણવામાં આપણને રસ હોય છે. તે માટે સંકેત રજૂ કરીએ.

\begin{defn}
  !!^a{formula}~$!A$ એ નિદર્શન $M = \tuple{W, R,
    V}$માં \emph{સત્ય} છે, જે $\mSat{M}{!A}$થી દર્શાવાય છે,
  ત્યારે અને માત્ર ત્યારે જ જ્યારે દરેક $w \in W$ માટે
  $\mSat{M}{!A}[w]$ હોય.
\end{defn}

\begin{prop}\ollabel{prop:truthfacts}
  \begin{enumerate}
  \item જો $\mSat{M}{!A}$ હોય તો $\mSat/{M}{\lnot !A}$ હોય,
    પણ તેનું વિપરીત વિધાન \emph{સાચું નથી}.
  \item જો $\mSat{M}{!A \lif !B}$ હોય તો $\mSat{M}{!A}$ હોય
    ત્યારે $\mSat{M}{!B}$ પણ હોય; પરંતુ તેનું વિપરીત વિધાન
    \emph{સાચું નથી}.
  \end{enumerate}
\end{prop}

\begin{proof}
  \begin{enumerate}
  \item જો $\mSat{M}{!A}$ હોય તો $!A$ એ $W$ના બધા જગતોમાં
    સત્ય છે. $W \neq \emptyset$ હોવાથી $\mSat{M}{\lnot!A}$
    શક્ય નથી; નહીંતર કોઈ એક જગતમાં $!A$ સત્ય પણ અને અસત્ય પણ
    હોવું પડે.

    બીજી બાજુ, જો $\mSat/{M}{\lnot !A}$ હોય તો $!A$ એ કોઈ
    $w \in W$માં સત્ય છે. તે પરથી \emph{દરેક}~$w \in W$ માટે
    $\mSat{M}{!A}[w]$ ફલિત થતું નથી. ઉદાહરણ તરીકે,
    \olref[rel]{fig:simple}ના નિદર્શનમાં $\mSat/{M}{\lnot p}$
    અને $\mSat/{M}{p}$ બંને છે.
  \item $\mSat{M}{!A \lif !B}$ અને $\mSat{M}{!A}$ માનો.
    $\mSat{M}{!B}$ બતાવવા માટે $w \in W$ કોઈપણ જગત લો.
    ત્યારે $\mSat{M}{!A \lif !B}[w]$ અને $\mSat{M}{!A}[w]$
    હોવાથી $\mSat{M}{!B}[w]$ છે. $w$ ગમે તે હતું, તેથી
    $\mSat{M}{!B}$.

    વિપરીત વિધાન ખોટું છે તે બતાવવા એવું નિદર્શન
    $\mModel{M}$ જોઈએ જેમાં $\mSat{M}{!A}$ હોય તો
    $\mSat{M}{!B}$ હોય, છતાં $\mSat/{M}{!A \lif !B}$ હોય.
    ફરી \olref[rel]{fig:simple}નું નિદર્શન લો:
    $\mSat/{M}{p}$ હોવાથી (પૂર્વપક્ષ ખોટો હોવાથી)
    $\mSat{M}{p}$ હોય તો $\mSat{M}{q}$ હોય એવું વિધાન સાચું છે.
    તેમ છતાં $\mSat/{M}{p \lif
      q}$ છે, કેમ કે~$w_1$માં $p$ સત્ય છે અને $q$ અસત્ય છે.
  \end{enumerate}
\end{proof}

\begin{prob}
  માત્ર $p_1$, $p_2$, $p_3$ને !!{propositional variable}s તરીકે
  ધરાવતી ભાષા માટે નીચેનું નિદર્શન $\mModel{M}$ વિચારો:
  \begin{center}
    \begin{tikzpicture}[modal]
      \node[world] (w1) [label={[align=right]left:\mTrue{p_1}\\\mFalse{p_2}\\\mFalse{p_3}}]
            {$w_1$} ; 
      \node[world] (w2) [label={[align=right]right:\mTrue{p_1}\\\mTrue{p_2}\\\mFalse{p_3}},
        below right=of w1]  {$w_2$}; 
      \node[world] (w3) [label={[align=right]right:\mTrue{p_1}\\\mTrue{p_2}\\\mTrue{p_3}},
        above right=of w2] {$w_3$};
      \draw[reflexive above] (w3) to (w3); 
      \draw[->] (w1) to (w2);
      \draw[->] (w2) to (w3);
      \draw[->] (w1) to (w3); 
    \end{tikzpicture}
  \end{center}
  નીચેના !!{formula}s અને સૂત્ર-ઢાંચા નિદર્શન
  $\mModel{M}$માં, એટલે કે $\mModel{M}$ના દરેક જગતમાં, સત્ય છે?
  સમજાવો.
  \begin{enumerate}
  \item $p\lif \Diamond p$ ($p$ પરમાણુ હોય ત્યારે);
  \item $!A\lif \Diamond !A$ ($!A$ કોઈપણ સૂત્ર હોય ત્યારે);
  \item $\Box p \lif p$ ($p$ પરમાણુ હોય ત્યારે);
  \item $\lnot p \lif \Diamond \Box p$ ($p$ પરમાણુ હોય ત્યારે);
  \item $\Diamond \Box !A$ ($!A$ કોઈપણ સૂત્ર હોય ત્યારે);
  \item $\Box \Diamond p$ ($p$ પરમાણુ હોય ત્યારે).
  \end{enumerate}
\end{prob}

\end{document}
